%% =====================================================================
%%  Unit 9 — Excuse me, where's the bank?  (Demander son chemin)
%%  Pistes : 9.01–9.04 dialogues · 9.05–9.06 vocabulaire · 9.07–9.08 grammaire
%%           9.09 prononciation · 9.10–9.13 exercices (textes en annexe)
%% =====================================================================

\newcommand{\otIllusUnIX}{\otIllusDuo{left/104/86/15/0.10,right/190/92/20/0.09,street/122/14/13/0.08,
  avenue/150/8/11/0.08,park/196/22/12/0.07,bank/100/50/12/0.07}{lea}{Excuse me, where's the bank?}{emma}{Turn left!}}

% Plan du quartier (exercices 3, 4, 8 et 11) — #1 : chemins de la solution
\newcommand{\otPlanVille}[1][]{%
  \begin{tikzpicture}[x=1mm,y=1mm,
    lieu/.style={rounded corners=1mm,font=\otSemi\fontsize{7.6}{9}\selectfont,align=center},
    rue/.style={font=\otMedium\fontsize{6.6}{8}\selectfont,text=otMuted}]
    \fill[otRule!45] (0,4) rectangle (150,96);
    % avenues (verticales) et rues (horizontales)
    \foreach \x in {10,75,140}{\fill[white] (\x-3,4) rectangle (\x+3,96);}
    \foreach \y in {10,50,90}{\fill[white] (0,\y-3) rectangle (150,\y+3);}
    \foreach \x/\n in {10/First Avenue,75/Second Avenue,140/Third Avenue}{%
      \node[rue,rotate=90] at (\x,30) {\enb{\n}};\node[rue,rotate=90] at (\x,70) {\enb{\n}};}
    \foreach \y/\n in {10/Main Street,50/Oak Street,90/River Street}{%
      \node[rue] at (42,\y) {\enb{\n}};\node[rue] at (107,\y) {\enb{\n}};}
    % lieux
    \fill[otGreen!35,rounded corners=1mm] (14,54) rectangle (71,86);
    \node[lieu,text=otGreenDark] at (42.5,70) {\enb{PARK}};
    \fill[otNavyTint!50!otNavy!18,rounded corners=1mm] (14,14) rectangle (42,46);
    \node[lieu] at (28,30) {\enb{bank}};
    \fill[otOrangeTint,rounded corners=1mm] (44,14) rectangle (71,46);
    \node[lieu] at (57.5,30) {\enb{post}\\\enb{office}};
    \fill[otVioletTint,rounded corners=1mm] (79,14) rectangle (107,46);
    \node[lieu] at (93,30) {\enb{super-}\\\enb{market}};
    \fill[otTealTint,rounded corners=1mm] (109,14) rectangle (136,46);
    \node[lieu] at (122.5,30) {\enb{pharmacy}};
    \fill[otRedTint,rounded corners=1mm] (79,54) rectangle (107,86);
    \node[lieu] at (93,70) {\enb{museum}};
    \fill[otNavyTint,rounded corners=1mm] (109,54) rectangle (136,86);
    \node[lieu] at (122.5,70) {\enb{hotel}};
    % départ
    \node[star,star points=5,star point ratio=2.2,fill=otOrange,inner sep=0pt,minimum size=4.2mm] (dep) at (75,0.5) {};
    \node[anchor=west,font=\otSemi\fontsize{7.6}{9}\selectfont,text=otOrangeDark] at (78,0.5) {Vous êtes ici};
    \draw[otOrangeDark,line width=0.6pt,-{Stealth[length=1.6mm]}] (71,0.5) -- (71,5.5) node[left,font=\otMedium\fontsize{6.4}{8}\selectfont] {N};
    #1
  \end{tikzpicture}}
\newcommand{\otChemin}[4]{% couleur, points, position du numéro, numéro
  \draw[#1,line width=1.4pt,line cap=round,line join=round,-{Stealth[length=2.6mm]},opacity=0.9] #2;
  \node[circle,fill=#1,text=white,inner sep=0pt,minimum size=4mm,font=\otBold\fontsize{7}{7}\selectfont] at #3 {#4};}

\unite[Excuse me,\\where's the bank?]{9}{Excuse me, where's the bank?}{Demander son chemin}{otIllusUnIX}

\situation{Épisode 9}{Samedi matin, à Manhattan. Léa doit envoyer un colis en France, mais où est
la poste ? Elle demande son chemin à un passant. L'après-midi, Emma lui donne rendez-vous au musée,
de l'autre côté de Central Park.}

\begin{objectifs}
\item demander poliment votre chemin : \en{Excuse me, where's…?} \en{How do I get to…?} ;
\item comprendre et donner des indications : tout droit, à gauche, à droite… ;
\item utiliser l'impératif : \en{Turn left}, \en{Don't take the bus} ;
\item nommer les lieux de la ville et situer un bâtiment sur un plan ;
\item distinguer \phon{ɪ} et \phon{iː} : \en{ship} et \en{sheep}.
\end{objectifs}

\auprogramme
  {Demander et indiquer un chemin}
  {L'impératif\par\vspace{0.6mm}\en{next to}, \en{opposite}, \en{on the corner of}}
  {La ville, les lieux, les transports : 29 expressions}
  {\phon{ɪ} ou \phon{iː} : \en{ship}, \en{sheep}}
  {Se repérer à Manhattan}

\semaine{casque/otNavy/Dialogues/Écouter et lire les deux scènes,
  bulle/otGreenDark/Vocabulaire/Apprendre les 29 expressions,
  blocs/otViolet/Grammaire/L'impératif et le plan,
  sons/otOrangeDark/Sons et culture/ship ou sheep et Manhattan,
  crayon/otNavy/Exercices/12 exercices corrigés,
  coche/otGreenDark/Bilan/L'essentiel et le test express}

\clearpage
%% ---------------------------------------------------------------------
\rubrique{ecouter}{Dans les rues de Manhattan}

\begin{dialogue}{1}{Excuse me, where's the post office?}{Samedi, 10 heures, à Manhattan. Léa a un
  colis à la main. Elle s'adresse à un passant.}{9.01}{9.02}
\dit{lea}{\hl{Excuse me}, where's the post office, please?\note{1}}{Excusez-moi, où est la poste, s'il vous plaît ?}
\dit{passant}{The post office? \hl{Go straight ahead}, then \hl{turn left} on Fifth Avenue.\note{2}}{La poste ? Allez tout droit, puis tournez à gauche sur la Cinquième Avenue.}
\dit{lea}{Straight ahead, then left on Fifth Avenue.}{Tout droit, puis à gauche sur la Cinquième Avenue.}
\dit{passant}{That's right. It's \hl{next to} a bank, \hl{opposite} a big park.}{C'est ça. Elle est à côté d'une banque, en face d'un grand parc.}
\dit{lea}{\hl{Is it far}?}{C'est loin ?}
\dit{passant}{No, it isn't far. It's \hl{a five-minute walk}.\note{3}}{Non, ce n'est pas loin. C'est à cinq minutes à pied.}
\dit{lea}{Thank you very much!}{Merci beaucoup !}
\dit{passant}{You're welcome. Oh, and \hl{don't take} the bus. The traffic is terrible!\note{4}}{Je vous en prie. Ah, et ne prenez pas le bus. La circulation est épouvantable !}
\end{dialogue}

\begin{comprendre}
\pointnote{1} \en{\textbf{Excuse me}} « Excusez-moi », pour aborder poliment un inconnu. On termine
  la question par \en{please}.
\pointnote{2} \en{\textbf{Go straight ahead, then turn left.}} Pour donner une indication, on emploie
  l'\textbf{impératif} : la base verbale seule, sans sujet. \en{turn left} = « tournez à gauche ».
\pointnote{3} \en{\textbf{a five-minute walk}} « à cinq minutes à pied ». Dans cette expression,
  \en{minute} reste au singulier.
\pointnote{4} \en{\textbf{Don't take the bus.}} L'impératif négatif se forme avec \en{Don't} + base
  verbale : « ne prenez pas le bus ».
\end{comprendre}

\begin{dialogue}{2}{How do I get to the museum?}{Samedi, 14 heures. Au téléphone, Emma explique à
  Léa comment la rejoindre au musée.}{9.03}{9.04}
\dit{emma}{Hi, Léa! Are you ready?}{Salut, Léa ! Tu es prête ?}
\dit{lea}{Yes! \hl{How do I get to} the museum?\note{5}}{Oui ! Comment je vais au musée ?}
\dit{emma}{\hl{Take the subway} from your apartment. \hl{Get off} at Central Park.}{Prends le métro depuis chez toi. Descends à Central Park.}
\dit{lea}{And then?}{Et ensuite ?}
\dit{emma}{\hl{Walk} across the park. The museum is \hl{on the right}, in front of a big fountain.\note{6}}{Traverse le parc à pied. Le musée est sur la droite, devant une grande fontaine.}
\dit{lea}{Across the park, on the right\ldots{} OK!}{Traverser le parc, sur la droite… D'accord !}
\dit{emma}{And \hl{don't be late}!}{Et ne sois pas en retard !}
\dit{lea}{Me? Late? Never! See you at three!}{Moi ? En retard ? Jamais ! À 15 heures !}
\end{dialogue}

\begin{comprendre}
\pointnote{5} \en{\textbf{How do I get to…?}} « Comment je vais à… ? », « comment se rendre à… ? ».
  C'est la question la plus utile pour se déplacer. \en{get off} = « descendre » (d'un bus, du métro).
\pointnote{6} \en{\textbf{on the right}} « sur la droite » ; \en{on the left} = « sur la gauche ».
  \en{across the park} = « à travers le parc ».
\end{comprendre}

%% ---------------------------------------------------------------------
\rubrique{mots}{Les mots de l'unité}

Voici les 29 expressions de l'unité. Retenez en particulier la différence entre \en{turn right}
(« tournez à droite ») et \en{go straight ahead} (« allez tout droit ») : en anglais, « droit » et
« droite » ne se ressemblent pas du tout !

\begin{vocabulaire}{Les lieux de la ville}{9.05}
\categorie{Les bâtiments}
\mot{a bank}{ə ˈbæŋk}{une banque}
\mot{a post office}{ə ˈpoʊst ˌɑːfɪs}{une poste}
\mot{a museum}{ə mjuːˈziːəm}{un musée}
\mot{a supermarket}{ə ˈsuːpɚˌmɑːrkɪt}{un supermarché}
\mot{a pharmacy}{ə ˈfɑːrməsi}{une pharmacie}
\mot{a hospital}{ə ˈhɑːspɪtəl}{un hôpital}
\mot{a hotel}{ə hoʊˈtel}{un hôtel}
\mot{a restaurant}{ə ˈrestərɑːnt}{un restaurant}
\categorie{Se déplacer}
\mot{a park}{ə ˈpɑːrk}{un parc}
\mot{a street}{ə ˈstriːt}{une rue}
\mot{an avenue}{ən ˈævənuː}{une avenue}
\mot{a bridge}{ə ˈbrɪdʒ}{un pont}
\mot{a station}{ə ˈsteɪʃən}{une gare}
\mot{a subway station}{ə ˈsʌbweɪ ˌsteɪʃən}{une station de métro}
\mot{a bus stop}{ə ˈbʌs ˌstɑːp}{un arrêt de bus}
\end{vocabulaire}

\begin{vocabulaire}{Demander et indiquer le chemin}{9.06}
\categorie{Demander}
\mot{Excuse me.}{ɪkˈskjuːz miː}{Excusez-moi.}
\mot{Where's the bank?}{ˌwerz ðə ˈbæŋk}{Où est la banque ?}
\mot{How do I get to the museum?}{ˌhaʊ duː aɪ ˈɡet tə ðə mjuːˈziːəm}{Comment je vais au musée ?}
\mot{Is it far?}{ɪz ɪt ˈfɑːr}{C'est loin ?}
\mot{It's near here.}{ɪts ˌnɪr ˈhɪr}{C'est près d'ici.}
\mot{It's a five-minute walk.}{ɪts ə ˌfaɪv ˌmɪnɪt ˈwɑːk}{C'est à cinq minutes à pied.}
\categorie{Indiquer}
\mot{go straight ahead}{ˌɡoʊ ˌstreɪt əˈhed}{aller tout droit}
\mot{turn left · turn right}{ˌtɝːn ˈleft · ˌtɝːn ˈraɪt}{tourner à gauche · à droite}
\mot{on the left · on the right}{ɑːn ðə ˈleft · ɑːn ðə ˈraɪt}{sur la gauche · sur la droite}
\mot{cross the street}{ˌkrɑːs ðə ˈstriːt}{traverser la rue}
\mot{take the subway}{ˌteɪk ðə ˈsʌbweɪ}{prendre le métro}
\mot{get off}{ˌɡet ˈɑːf}{descendre (du métro, du bus)}
\mot{on the corner of}{ɑːn ðə ˈkɔːrnɚ əv}{à l'angle de}
\mot{across from}{əˈkrɑːs frəm}{en face de \registre{américain}}
\end{vocabulaire}

%% ---------------------------------------------------------------------
\rubrique{grammaire}{L'impératif et les indications}

\otIntertitre[otViolet]{1 · L'impératif}
\begin{paire}
\begin{regle}[Affirmatif : la base verbale seule]
\exen{Go straight ahead.}{Allez tout droit.}
\exen{Turn left.}{Tournez à gauche.}
\exen{Take the subway.}{Prends le métro.}
\exen{Please, sit down.}{Asseyez-vous, je vous prie.}
\end{regle}
\begin{regle}[Négatif : \en{Don't} + base verbale]
\exen{Don't take the bus.}{Ne prenez pas le bus.}
\exen{Don't be late!}{Ne sois pas en retard !}
\exen{Don't worry.}{Ne t'inquiète pas.}
\exen{Don't cross here.}{Ne traversez pas ici.}
\end{regle}
\end{paire}
{\fontsize{10}{13}\selectfont L'impératif est le même pour « tu » et pour « vous ». Pour adoucir une
demande, ajoutez \en{please}, au début ou à la fin.\par}

\blocaudio{9.07}{Écoutez et répétez}{Six indications à l'impératif, affirmatives puis négatives.}

\otIntertitre[otViolet]{2 · Demander son chemin}
\begin{center}\otTableau
\begin{tabular}{@{}>{\enOn\bfseries}l<{\enOff}@{\hspace{8mm}}l@{}}
Excuse me, where's the post office?    & Excusez-moi, où est la poste ?\\
How do I get to the station?           & Comment je vais à la gare ?\\
Is there a bank near here?             & Y a-t-il une banque près d'ici ?\\
Is it far? — No, it's near here.       & C'est loin ? — Non, c'est près d'ici.\\\bottomrule
\end{tabular}\end{center}

\blocaudio{9.08}{Écoutez et répétez}{Quatre questions pour demander son chemin, et leur réponse.}

\begin{attention}[Les pièges des francophones]
\fauxjuste{Turn at left.}{Turn left.}
\fauxjuste{Go right ahead.}{Go straight ahead.}
\fauxjuste{Don't to take the bus.}{Don't take the bus.}
\fauxjuste{Where is the bank, please, Sir?}{Excuse me, where's the bank?}
\par\vspace{1mm}{\fontsize{10}{12.6}\selectfont « Tout droit » se dit \en{straight ahead} ;
\en{right} veut dire « à droite ». Ne les confondez pas : vous enverriez votre interlocuteur dans
la mauvaise direction !\par}
\end{attention}

%% ---------------------------------------------------------------------
\rubrique[ship ou sheep]{sons}{\phon{ɪ} ou \phon{iː} : \en{ship} ou \en{sheep}}

L'anglais a deux « i » : le \phon{iː} \textbf{long et tendu} de \en{sheep} (lèvres étirées, comme
pour sourire) et le \phon{ɪ} \textbf{bref et relâché} de \en{ship} (presque un « é »). Le français
n'a que le premier : les francophones prononcent souvent \en{live} comme \en{leave}, et \en{ship}
comme \en{sheep}. Or ces mots n'ont pas du tout le même sens !

\begin{paires}{9.09}{Le premier mot de chaque paire a le son bref \phon{ɪ}, le second le son long
\phon{iː}. Écoutez, puis répétez en exagérant la différence de longueur.}
\sonpaire{to live}{tə lɪv}{habiter, vivre}{to leave}{tə liːv}{partir}
\sonpaire{rich}{rɪtʃ}{riche}{reach}{riːtʃ}{atteindre}
\sonpaire{fit}{fɪt}{en forme}{feet}{fiːt}{les pieds}
\sonpaire{ship}{ʃɪp}{un navire}{sheep}{ʃiːp}{un mouton}
\sonpaire{fill}{fɪl}{remplir}{feel}{fiːl}{sentir, ressentir}
\sonpaire{will}{wɪl}{la volonté}{wheel}{wiːl}{une roue}
\sonpaire{list}{lɪst}{une liste}{least}{liːst}{le moins}
\sonpaire{slip}{slɪp}{glisser}{sleep}{sliːp}{dormir}
\end{paires}

\begin{astuce}[Et l'intonation ?]
Dans une question fermée (réponse \en{yes} ou \en{no}), la voix \textbf{monte} à la fin :
\en{Is it far?} ↗. Dans une question avec un mot interrogatif, elle \textbf{descend} :
\en{Where's the bank?} ↘. Écoutez les dialogues en y prêtant attention.
\end{astuce}

%% ---------------------------------------------------------------------
\rubrique{culture}{Se repérer à Manhattan}

\begin{culture}[Un quadrillage presque parfait]
Au nord de \en{Greenwich Village}, Manhattan est un quadrillage : les \en{avenues} vont du nord au
sud, les \en{streets} d'est en ouest, et la plupart sont numérotées. On mesure les distances en
\en{blocks} (pâtés de maisons) : \en{It's three blocks from here}. Aller vers le nord, c'est aller
\en{uptown} ; vers le sud, \en{downtown}.
\end{culture}

\begin{paire}
\begin{culture}[Le métro]
Le métro new-yorkais (\en{the subway}) circule 24 heures sur 24, tous les jours. Les lignes portent un
numéro ou une lettre (\en{the A train}, \en{the 6}). Aujourd'hui, on paie le plus souvent sans contact,
avec sa carte bancaire ou son téléphone.
\end{culture}
\begin{culture}[Demander poliment]
Commencez toujours par \en{Excuse me}, et remerciez : \en{Thank you so much!} Les New-Yorkais ont la
réputation d'être pressés, mais ils indiquent volontiers leur chemin aux visiteurs.
\end{culture}
\end{paire}

%% ---------------------------------------------------------------------
\rubrique{pratique}{Exercices}

\exercice[ecrire,niveau=1]{Où va-t-on ?}{Complétez avec un lieu de la ville.}
\begin{anglais}
\begin{questions}
\item You send a letter at the \trou[40mm]{post office}.
\item You see paintings at the \trou[40mm]{museum}.
\item You buy food at the \trou[40mm]{supermarket}.
\item You sleep at a \trou[40mm]{hotel}.
\item You take a train at the \trou[40mm]{station}.
\item You buy medicine at the \trou[40mm]{pharmacy}.
\end{questions}
\par\vspace{1mm}{\fontsize{8.6}{10}\selectfont\color{otMuted}\foreignlanguage{french}{\textit{send a letter} : envoyer une lettre · \textit{paintings} : des tableaux · \textit{medicine} : des médicaments}}
\end{anglais}

\exercice[ecrire,niveau=1]{À l'impératif}{Traduisez ces indications.}
\begin{questions}[2][french]
\item Tournez à gauche. → \enb{\trou[40mm]{Turn left.}}
\item Allez tout droit. → \enb{\trou[40mm]{Go straight ahead.}}
\item Ne prenez pas le bus. → \enb{\trou[40mm]{Don't take the bus.}}
\item Prends le métro. → \enb{\trou[40mm]{Take the subway.}}
\item Traversez la rue. → \enb{\trou[40mm]{Cross the street.}}
\item Ne sois pas en retard ! → \enb{\trou[36mm]{Don't be late!}}
\end{questions}

\exercice[lire,niveau=2]{Suivez le guide}{Partez de l'étoile, en regardant vers le nord. Lisez chaque
itinéraire et écrivez le lieu où vous arrivez.}
\begin{center}\otPlanVille\end{center}
\begin{questions}
\item Go straight ahead on Second Avenue. Turn left on Oak Street. It's on the right. \hfill \enb{\trou[30mm]{the park}}
\item Go straight ahead. Cross Main Street. It's on your right. \hfill \enb{\trou[30mm]{the supermarket}}
\item Go straight ahead. Turn right on Oak Street. It's on the left, next to the hotel. \hfill \enb{\trou[30mm]{the museum}}
\item Go straight ahead, then turn left on Main Street. It's on the right, next to the bank. \hfill \enb{\trou[30mm]{the post office}}
\end{questions}

\exercice[lire,ecrire,niveau=2]{Sur le plan}{Regardez le plan de l'exercice 3 et complétez avec
\en{next to}, \en{opposite}, \en{between} ou \en{on the corner of}.}
\begin{questions}
\item The bank is \trou[30mm]{next to} the post office.
\item The museum is \trou[30mm]{opposite} the park.
\item The post office is \trou[30mm]{between} the bank and Second Avenue.
\item The hotel is \trou[30mm]{on the corner of} Third Avenue and Oak Street.
\item The supermarket is \trou[30mm]{opposite} the post office.
\end{questions}

\exercice[lire,niveau=1]{Le bon ordre}{Remettez ce dialogue dans l'ordre : écrivez les lettres.}
\begin{anglais}
\begin{questions}[1][english]
\item[a.] No, it's a two-minute walk.
\item[b.] Excuse me, is there a pharmacy near here?
\item[c.] Thank you very much!
\item[d.] Yes, there is. Go straight ahead and turn right. It's opposite the bank.
\item[e.] You're welcome.
\item[f.] Is it far?
\end{questions}
\end{anglais}
\par\vspace{2mm}\noindent{\otSemi\fontsize{10}{12}\selectfont Ordre :}\hspace{4mm}%
\trou[9mm]{b}\hspace{3mm}\trou[9mm]{d}\hspace{3mm}\trou[9mm]{f}\hspace{3mm}\trou[9mm]{a}\hspace{3mm}\trou[9mm]{c}\hspace{3mm}\trou[9mm]{e}

\exercice[ecrire,niveau=2]{Demandez poliment}{Écrivez la question polie qui correspond à la situation.}
\begin{questions}[1][french]
\item Vous cherchez la gare. \enb{(where)} \hfill \enb{\trou[78mm]{Excuse me, where's the station?}}
\item Vous voulez aller au musée. \enb{(how)} \hfill \enb{\trou[78mm]{Excuse me, how do I get to the museum?}}
\item Vous cherchez une banque dans le quartier. \enb{(is there)} \hfill \enb{\trou[78mm]{Excuse me, is there a bank near here?}}
\item Vous voulez savoir si c'est loin. \hfill \enb{\trou[78mm]{Is it far?}}
\end{questions}

\exercice[ecouter,niveau=1,piste=9.10]{\phon{ɪ} ou \phon{iː} ?}{Écoutez. Pour chaque numéro, la voix dit
un seul des deux mots, deux fois. Cochez le mot entendu.}
\begin{questions}[4]
\item \choix{live}\quad\choix{leave}
\item \choix{rich}\quad\choix{reach}
\item \choix{fit}\quad\choix{feet}
\item \choix{ship}\quad\choix{sheep}
\item \choix{fill}\quad\choix{feel}
\item \choix{will}\quad\choix{wheel}
\item \choix{list}\quad\choix{least}
\item \choix{slip}\quad\choix{sleep}
\end{questions}

\exercice[ecouter,lire,niveau=2,piste=9.11]{Où vont-ils ?}{Trois personnes partent de l'étoile du plan
(exercice 3). Écoutez les indications et écrivez le lieu où elles arrivent.}
\begin{questions}[3]
\item \enb{\trou[30mm]{the hotel}}
\item \enb{\trou[30mm]{the pharmacy}}
\item \enb{\trou[30mm]{the bank}}
\end{questions}

\exercice[ecouter,ecrire,niveau=2,piste=9.12]{Dictée à trous}{Un touriste demande son chemin à Jack.
Écoutez et complétez.}
\begin{anglais}
\noindent\begin{tabular}{@{}l@{\hspace{6mm}}l@{}}
\otNomDit{passant} & \trou{Excuse} me, how do I get to the station?\\[1.8mm]
\otNomDit{jack}    & Go \trou{straight} ahead, then turn \trou{right}.\\[1.8mm]
\otNomDit{passant} & Is it \trou{far}?\\[1.8mm]
\otNomDit{jack}    & No, it's a ten-minute \trou{walk}. It's \trou{opposite} the museum.\\[1.8mm]
\otNomDit{passant} & Can I take the bus?\\[1.8mm]
\otNomDit{jack}    & Yes, but \trou{don't} take the number 4. \trou{Take} the number 2.\\
\end{tabular}
\end{anglais}

\exercice[ecrire,niveau=2]{Traduisez}{Traduisez en anglais.}
\begin{questions}[1][french]
\item Excusez-moi, où est la gare, s'il vous plaît ?\lignes{1}
\item Allez tout droit, puis tournez à droite.\lignes{1}
\item C'est loin ? — Non, c'est à cinq minutes à pied.\lignes{1}
\item La banque est à côté de la poste.\lignes{1}
\item Ne prenez pas le métro !\lignes{1}
\item Descendez à Central Park.\lignes{1}
\item Y a-t-il un supermarché près d'ici ?\lignes{1}
\item Le musée est sur la gauche, en face du parc.\lignes{1}
\end{questions}

\exercice[ecrire,niveau=3]{Votre itinéraire}{Sur le plan de l'exercice 3, écrivez l'itinéraire de
l'étoile jusqu'à la pharmacie, en trois ou quatre phrases. Puis écrivez le chemin de chez vous jusqu'à
votre boulangerie (\en{bakery}) préférée.}
\lignes{6}

\exercice[ecouter,parler,niveau=3,piste=9.13]{À vous de parler !}{Vous êtes à l'étoile du plan.
Une touriste vous demande son chemin. Répondez \textbf{à voix haute} pendant chaque pause.}
\begin{center}\otTableau
\begin{tabular}{@{}>{\otSemi\fontsize{8.6}{10}\selectfont\color{otNavy}}l@{\hspace{3mm}}>{\enOn}p{70mm}<{\enOff}@{\hspace{6mm}}p{78mm}@{}}
 & \otTH{La touriste dit\ldots} & \otTH{Vous\ldots}\\\midrule
1. & Excuse me, where's the museum, please? & indiquez le chemin (tout droit, à droite…).\\
2. & Is it far? & répondez que non : cinq minutes à pied.\\
3. & And is there a pharmacy near here? & répondez et indiquez où elle est.\\
4. & Can I take the bus? & conseillez-lui de ne pas prendre le bus.\\
5. & Thank you so much! & répondez poliment.\\\bottomrule
\end{tabular}\end{center}

%% ---------------------------------------------------------------------
\rubrique{bilan}{Faites le point}

\begin{essentiel}
\begin{minipage}[t]{0.57\linewidth}
\otEssTitre{Les phrases clés}
\renewcommand{\arraystretch}{1.32}
\begin{tabular}{@{}>{\enOn\otSemi\fontsize{9.4}{11}\selectfont}l<{\enOff}@{\hspace{3mm}}>{\fontsize{9.2}{11}\selectfont\color{white!78!otNavy}}l@{}}
Where's the bank?              & Où est la banque ?\\
How do I get to the museum?    & Pour aller au musée ?\\
Go straight ahead.             & Allez tout droit.\\
Turn left. / Turn right.       & Tournez à gauche / droite.\\
It's on the right.             & C'est sur la droite.\\
Is it far? — It's near here.   & C'est loin ? — C'est près.\\
\end{tabular}
\end{minipage}\hfill
\begin{minipage}[t]{0.38\linewidth}
\otEssTitre{L'impératif}
{\enb{\fontsize{9.4}{12}\selectfont \textbf{Take} the subway. · \textbf{Don't} take the bus.}\par}
\otEssTitre{Situer}
{\enb{\fontsize{9.2}{11.6}\selectfont next to · opposite · between · on the corner of · on the left / right}\par}
\otEssTitre{Prononciation}
{\fontsize{9.2}{11.6}\selectfont \phon{ɪ} bref : \en{live}, \en{ship} · \phon{iː} long : \en{leave}, \en{sheep}\par}
\end{minipage}
\end{essentiel}

\begin{jesais}
\item aborder poliment un passant et demander mon chemin ;
\item comprendre des indications simples et les suivre sur un plan ;
\item donner un itinéraire à l'impératif, affirmatif et négatif ;
\item nommer une quinzaine de lieux de la ville ;
\item situer un bâtiment : \en{next to}, \en{opposite}, \en{between}, \en{on the corner of} ;
\item distinguer \phon{ɪ} et \phon{iː} à l'écoute et à l'oral.
\end{jesais}

\clearpage
\begin{testexpress}
\qcm[a]{« Excusez-moi » (pour aborder quelqu'un) se dit :}{\en{Excuse me.}}{\en{Sorry me.}}{\en{Pardon me you.}}
\qcm[c]{« Allez tout droit » se dit :}{\en{Go right.}}{\en{Go on the right.}}{\en{Go straight ahead.}}
\qcm[b]{« Tournez à gauche » se dit :}{\en{Turn at left.}}{\en{Turn left.}}{\en{Turn on left.}}
\qcm[b]{« Ne prenez pas le bus » se dit :}{\en{No take the bus.}}{\en{Don't take the bus.}}{\en{Don't to take the bus.}}
\qcm[a]{« Comment je vais à la gare ? » se dit :}{\en{How do I get to the station?}}{\en{How I go the station?}}{\en{Where I go to the station?}}
\qcm[c]{« en face du parc » se dit :}{\en{next to the park}}{\en{behind the park}}{\en{opposite the park}}
\qcm[a]{« C'est à cinq minutes à pied » se dit :}{\en{It's a five-minute walk.}}{\en{It's five minutes of foot.}}{\en{It's a five-minutes walk.}}
\qcm[b]{\en{get off} veut dire :}{monter}{descendre (du métro)}{attendre}
\qcm[c]{Quel mot a le son long \phon{iː} ?}{\en{ship}}{\en{live}}{\en{sheep}}
\qcm[a]{Aller vers le nord de Manhattan, c'est aller :}{\en{uptown}}{\en{downtown}}{\en{crosstown}}
\end{testexpress}

\noindent{\fontsize{9.6}{12}\selectfont\color{otMuted}Niveau visé (CECRL, niveau A1) : « peut
demander son chemin et comprendre des indications simples ». Les corrigés de l'unité sont
page~\pageref{cor-u9-1} ; les textes des pistes d'exercices, page~\pageref{trans-9.10}.\par}
